\Needspace{6\baselineskip}
\begin{center}\hypertarget{native-251-article179}{}\pdfbookmark[2]{179.}{article-179}179.\end{center}
We shall now develop the consequences of this law. It is clear that for the product $\Omega=h^\mu\pi^{-\frac{1}{2}\mu}e^{-hh(vv+v'v'+v''v''+\mathord{.}\,\mathord{.})}$ to be a maximum, the sum $vv+v'v'+v''v''+\mathord{.}\,\mathord{.}$ must be a minimum. \emph{The most probable set of values for the unknowns $p$, $q$, $r$, $s$ etc. is therefore the one for which the squares of the differences between the observed and computed values of the functions $V$, $V'$, $V''$ etc. give the least sum,} provided that the same degree of precision is to be assumed in all observations.

This principle, which finds frequent use in all applications of mathematics to natural philosophy, should have the force of an axiom everywhere, with the same right as the arithmetic mean of several observed values of a single quantity is taken as its most probable value.

The principle can now readily be extended to observations of \emph{unequal} precision. Suppose that the measure of precision of the observations by which $V=M$, $V'=M'$, $V''=M''$ etc. were found is expressed respectively by $h$, $h'$, $h''$ etc., that is, that errors inversely proportional to these quantities could have occurred equally readily in those observations. This is plainly the same as if observations of equal precision (whose measure is $=1$) had yielded the values of the functions $hV$, $h'V'$, $h''V''$ etc. directly as $=hM$, $h'M'$, $h''M''$ etc. \SourceLeaf{252}The most probable set of values for the quantities $p$, $q$, $r$, $s$ etc. is therefore the one for which the sum $hhvv+h'h'v'v'+h''h''v''v''+\text{etc.}$, that is, \emph{the sum of the squares of the differences between the values actually observed and the computed values, each difference multiplied by the number that measures the degree of precision, is a minimum.} Accordingly, the functions $V$, $V'$, $V''$ etc. need not even refer to homogeneous quantities. They can also represent heterogeneous quantities (for example, seconds of arc and seconds of time), provided that we can estimate the ratio of the errors that could have occurred equally readily in the respective observations.
